% Generated draft: requires mathematical and source review.
\documentclass[11pt]{article}
\usepackage[margin=25mm]{geometry}
\usepackage{fontspec}
\setmainfont{Latin Modern Roman}
\usepackage{amsmath,amssymb,mathtools,unicode-math}
\setmathfont{Latin Modern Math}
\usepackage{enumitem,needspace,xurl,hyperref}
\setlength{\parindent}{0pt}
\setlength{\parskip}{5pt}
\setlength{\emergencystretch}{3em}
\newcommand{\N}{\mathbb{N}}
\newcommand{\Z}{\mathbb{Z}}
\newcommand{\Q}{\mathbb{Q}}
\newcommand{\R}{\mathbb{R}}
\newcommand{\C}{\mathbb{C}}
\newcommand{\F}{\mathbb{F}}
\newcommand{\A}{\mathbb{A}}
\newcommand{\PP}{\mathbb{P}}
\newcommand{\E}{\mathbb{E}}
\newcommand{\eps}{\varepsilon}
\newcommand{\dd}{\,\mathrm{d}}
\begin{document}
\section*{UnsolvedMath ID: 30004514}
\textbf{Source problem code:} catalogue entry
\subsection{Definitions and mathematical statement}
For an ambient dimension \(n\geq 1\), use the standard convention that a convex body \(K\subset\R^n\) is a compact convex set with nonempty interior. Let \(\lambda_n\) denote \(n\)-dimensional Lebesgue measure and define the barycenter of \(K\) by
\[
 c_K=\frac{1}{\lambda_n(K)}\int_K x\,\dd\lambda_n(x).
\]
An affine hyperplane is a set
\[
 h(u,t)=\{x\in\R^n:\langle u,x\rangle=t\},
 \qquad u\in\R^n\setminus\{0\},\quad t\in\R.
\]
For a hyperplane \(h\) containing \(c_K\), equip \(h\) with its induced \((n-1)\)-dimensional Euclidean measure \(\lambda_h\). Since \(c_K\) is an interior point of \(K\), its section has positive finite measure. Define
\[
 c_{K\cap h}=\frac{1}{\lambda_h(K\cap h)}
 \int_{K\cap h}x\,\dd\lambda_h(x).
\]
In dimension \(n=1\), the induced zero-dimensional measure is counting measure. Call \(h\) a barycentric hyperplane of \(K\) through \(c_K\) if
\[
 c_K\in h\qquad\text{and}\qquad c_{K\cap h}=c_K.
\]
Hyperplanes are counted as distinct subsets of \(\R^n\), without orientations or multiplicities.

\textbf{Question as supplied.} For every convex body \(K\subset\R^n\), must there exist pairwise distinct affine hyperplanes \(h_1,\ldots,h_{n+1}\) such that
\[
 c_K\in h_i\quad\text{and}\quad c_{K\cap h_i}=c_K
 \qquad (i=1,\ldots,n+1)?
\]
\textbf{Unresolved dimensional qualification.} The supplied record gives no lower bound on \(n\). The intended range of dimensions, particularly the treatment or exclusion of \(n=1\), requires clarification; no restriction such as \(n\geq2\) is silently added here.

The supplied dated curation labels this question open and states that the proposed \(n+1\) assertion was reopened following a proof error, with lower bounds such as four cuts in dimensions at least three reported as known. These are supplied literature assessments, not independently verified conclusions of this draft.
\subsection{Short English statement}
Does every convex body in \(\R^n\) have at least \(n+1\) distinct hyperplanes through its centroid whose sections also have that centroid? Centroids use uniform Euclidean volume in the appropriate dimension. The source leaves the intended range of \(n\) unspecified. Its open-status label and account of a proof error are supplied dated curation requiring independent review.
\subsection{Sources}
\begin{itemize}
\item Supplied local UnsolvedMath registry record, permanent ID \(30004514\), exact source problem code \texttt{catalogue entry}, in the Oberwolfach Reports Open Problems collection. Title: \emph{Barycentric Hyperplanes Through Convex Body Centroids}. No individual proposer is identified; the supplied proposed year is 2020.
\item Supplied report citation: \emph{Discrete Geometry (2021)}, Oberwolfach Reports, pp.\ 1469--1528, DOI \texttt{10.4171/owr/2020/30}. \url{https://doi.org/10.4171/owr/2020/30}. The record describes the report as documenting workshop discussions and open problems; it supplies no exact page or passage for this question.
\item Patakova--Tancer--Wagner, \emph{Barycentric Cuts}, as attributed in the supplied literature sources. \url{https://doi.org/10.4230/LIPIcs.SoCG.2020.62}.
\item Patakova--Tancer--Wagner, preprint, as attributed in the supplied literature sources. \url{https://arxiv.org/abs/2003.13536}.
\item Literature assessment dated exactly \texttt{2026-08-22}: the supplied curation labels the problem open, reports reopening of the proposed \(n+1\) statement after a proof error, and reports lower bounds such as four cuts in dimensions at least three. This is supplied dated curation derived from the OWR record, not independently verified source consultation or a current literature assessment.
\item The registry labels the statement \texttt{corrected\_verified}, explaining that only formula delimiters, notation, grammar, and typography were repaired while mathematical content and scope were preserved. This supplied editorial label does not establish mathematical correctness or independent verification.
\end{itemize}
\end{document}
